\textbf{ Fichero general de prueba }

  Este fichero no contiene informacion ordenada o con una finalidad.
  \textit { Solo lo uso para testar nuevas features. } \newline

  Es por ello que tiene espacios, no mantiene una estructura clara y tampoco
  sigue un patron particular, inclusive, agregare de forma intencional malas
  practicas, con la intencion de llevar al limite este parser.

  \chapter{ Inicio }

  \section{ Primeras ideas }

  Este proyecto nacio de querer entender como se crea un lenguaje, como se define
  y se consume su sintaxis. De ahi surgen preguntas naturales como ¿por que
  hay diferencia de velocidad entre lenguajes?

  La respuesta directa es que depende de si es compilado o interpretado, pero
  la pregunta real es otra: ¿por que optaron por ese diseño y como cambia el
  funcionamiento del compilador?

  La primera idea fue parsear con \note[expresiones regulares]{Patrones para identificar
  y capturar texto} buscando estructuras tipo \backslash command\{.*\}\{.*(Contenido)\}.
  Al inicio parece util, pero al anidar comandos el regex falla: el \} del comando
  interior tambien cierra al superior. Se puede llevar un contador de \textbf{profundidad},
  pero la expresion crece rapido y en entornos se pone peor pues ya se necesitan
  funciones extra para extraer el contenido hasta el cierre.

  \section{ Otras alternativas }

  Probe con \note[Lark]{Herramienta para parseo y analisis gramatico} pero las 
  anidaciones tampoco funcionaban de forma clara.
  La siguiente idea fue trabajar caracter a caracter: acumular hasta encontrar
  un caracter especial y reaccionar segun el contexto.

  El problema es que el contexto complica todo. Un \{ en modo matematico es un
  literal, no una accion. Eso genera demasiadas validaciones. Es como resolver
  2+2 mediante union de conjuntos: mucho control para algo que podia ser simple.

  \section{ Pausa }

  Con esos intentos sin resultado claro, detuve el proyecto. No habia una version
  \note[estable]{Una que compilara al menos el documento mas basico} lista.
  Algunos estudiantes me contactaron para apoyar, les agradeci pero sin un
  parser estable la estructura podia cambiar en meses y

  \alert{ pudiera pasar que este proyecto se abandonara. }
  
  Meses despues me someti a una cirugia, nada complicado pero que me brindo unos
  dias de reposo y fue durante este tiempo que se me ocurrio continuar con este
  proyecto, previamente abandonado.

  \chapter{ Tokens }

  Experimentando con algunos modelos de lenguaje, se me ocurrio replicar el uso
  de tokens: usar regex pero para crear una \textbf{tabla de tokens}. No importa
  el contenido, solo el tipo, si es un comando, texto, entorno, etc. Ademas note
  que el arbol resultante no necesita pensar en como renderizar sino en como
  organizar la salida, el navegador se encarga del resto.

  Los primeros comandos integrados fueron una prueba: \backslash section,
  \backslash\backslash y estilos de texto, para verificar que la tokenizacion
  funcionaba.

  \section{ Estructura de archivos }

  Cuando el tokenizador crecio, separar en archivos de una sola responsabilidad
  fue necesario. Seguir el flujo en un unico archivo se habia vuelto dificil,
  y el parseo asincrono tenia dos puntos de ejecucion por step. Se elimino.

  \section{ Entornos }

  Con esa base mas limpia, los primeros entornos fueron \texttt{enumerate}
  e \texttt{itemize}. En retrospectiva no fue la mejor eleccion porque estos
  entornos usan \backslash item como separador, un \textit{cierre implicito}
  sin marca final. Eso complica la deteccion, pero si funciona bien la
  misma logica sirve para otros elementos similares.

  \begin{itemize}
      \item Test de lista
      \item Prueba con inline math $\int 2x^{3y}$
      \item Cierre implicito del anterior
  \end{itemize}

  \chapter{ Contextos }

  \section{ Ecuaciones }

  Al parsear, todo se categoriza sin conocer el orden ni la posicion. Una ecuacion
  llega a trozos y dentro de un entorno de ecuacion no se permiten otros entornos.
  La solucion fue simple: solo leer, no procesar. Un parametro
  \note[raw]{Modo donde el contenido llega literal al arbol sin procesarse}
  indica que el contenido llega tal cual al arbol.

  \begin{equation}
      E = mc^2
  \end{equation}

  \section{ Comentarios }

  Algo que pase por alto. En LaTeX se usan mucho, a veces para quitar elementos,
  otras como referencias. Agregarlos es directo: luego de \%, todo hasta el
  salto de linea se ignora. Al no ser un comando ni entorno, se puede eliminar
  desde la categorizacion.

  \section{ Notas }

  El primer elemento inline con logica propia fue
  \note[note]{\backslash note crea un tooltip con contenido emergente}.
  Su implementacion fue un poco complicada ya que el elemento debia mantener
  una posicion relativa a su ejecutor, pero al mismo tiempo debia tener una
  posicion absoluta en el documento.

  Con las notas surgio la necesidad de restringir el contenido permitido
  dentro de cada elemento, ie, que pasa si intentas poner una
  \note[ecuacion]{Elemento demasiado grande para caber en un tooltip}
  dentro de una nota. Por eso existe \textit{inline\_only}, para limitar
  que se puede escribir segun el contexto.

  \section{ Imagenes }

  Otro elemento evidente para un sitio web. El contenido de estos comandos
  no debe parsearse de forma recursiva, llega plano al arbol como un raw.
  Se agrego el entorno \texttt{figure} con \backslash includegraphics
  y \backslash caption.

  \begin{figure}
    \includegraphics[80]{img/topologia.png}
    \caption{Prueba de imagen}
  \end{figure}

  Con el parser mas estable ya se puede pensar no solo en que funcione,
  sino en establecer buenas practicas desde la lectura del fichero tex.
  La estructura recomendada del folder es:

\begin{verbatim}
book
 |-- main.tex
 |-- assets
       |-- pictures
              |-- 1.png
              |-- 2.jpg
\end{verbatim}

  y para leerla se debe usar su ruta relativa, es decir \texttt{ assets/pictures/1.png }

  \chapter{ Custom components }

  \section{ Bloques }

  Agrupar elementos en un
  \note[bloque]{Entorno para mostrar distintas versiones de un contenido con botones}
  nos permite indicar directamente contenidos que estan relaionados pero que de alguna
  forma, tienen una relacion aun mas estrecha que el resto del contexto.
  La primera idea fue un comando \backslash block\{c1\}\{c2\}\{...\}, pero la
  cantidad de subbloques es variable y eso complica el registro y la deteccion del cierre.
  Usar items resulto ser la solucion mas practica.

  \begin{block}
      \item Son utiles para mostrar distintas demostraciones del mismo elemento.
          Por ejemplo, en Topologia se puede abordar desde metricas o desde el
          concepto de \note[bolas]{Un subconjunto de $R^n$}. Aunque las bolas
          aparecen como consecuencia de tener una metrica, tambien son un
          concepto independiente.
      \item El primer intento con \backslash block\{c1\}\{c2\}\{...\} falla
          porque la cantidad variable de subbloques complica crear el registro
          y determinar el cierre.
      \item Usar items e indentacion resulta mucho mas simple.
  \end{block}

  \section{ Referencias }

  Cada \backslash\note[label]{Genera una entrada en un diccionario con un contador automatico}
  genera una entrada en \texttt{label:\{id, index\}}.
  Esto tambien cambio el output: ya no es un array sino un diccionario, pues
  las referencias viven como hermanos del arbol.

  \begin{equation}
      E = mc^2
  \end{equation}
  \label{eq:einstein}

  Se citan con \backslash ref\{id\}, ej, \ref{eq:einstein}.

  \section{ Componentes de la libreria }

  Los custom components de la libreria existian antes del parser. Integrarlos
  requirio retrocompatibilidad: se convirtieron de comandos inline a entornos
  con key-value props. Asi siguen la misma sintaxis que el parser ya conoce y
  el escritor no necesita aprender nada nuevo.

  \begin{axiom}
    Los custom components dejaron de ser commands inline para ser environments con key-value props.
  \end{axiom}

  \begin{axiom}[title=Arquimedes, label=ax-arq]
    Para todo $x \in \N$ existe $y \in \N$ tal que $y > x$.
  \end{axiom}

  \begin{theorem}[Bolzano-Weierstrass]
    Toda sucesion acotada en $\R^n$ tiene una subsucesion convergente.
  \end{theorem}

  \begin{definition}[title=Espacio metrico, label=def-met]
    Un espacio metrico es un par $(X, d)$ donde ...
  \end{definition}

  \begin{proof}
    Sea $\epsilon > 0$ ...
    \qed
  \end{proof}

  \begin{exercise}[label=ej-1]
    Demuestre que ...
  \end{exercise}

  \chapter{ Capitulos y estilos }

  \section{ Capitulos }

  Sin capitulos y secciones todo el contenido vive en un solo nivel: no hay
  tabla de contenido ni anclas de navegacion, algo esencial para documentos largos.

  \section{ Estilos }

  Con el parser estable, muchos documentos aun se rompen por comandos comunes
  que no estan registrados. Pero integrarlos no es complicado, \textsc{solo hay
  que registrarlos} y el parser ya sabe como tratarlos. Se agregaron estilos de
  texto y urls [ \url{https://cris0501.github.io/library/doc} ].

  Probablemente integre soporte para grupos y ahi pause el \textbf{desarrollo},
  pues ya seria momento de pasar algunas notas antes de agregar mas features.
  Solo el parser requiere mucho tiempo, y crear un \texttt{expansor de macros}
  como LaTeX es trabajo largo.

  Hoy cambie el contenido de este \textit{fichero}: antes no tenia sentido,
  solo probaba features. Ahora que el parser lee bien la estructura de un texto,
  el contenido puede tener coherencia.

